% titel: Blatt 9, A4: Osmotischer Druck
% skript: Mortimer Kap. 13 (Osmose)

% Teil I: Werkzeuge, kein Lösungsweg.

Durch eine semipermeable Membran treibt gelöstes Material Wasser
an, bis ein Gegendruck dies ausgleicht --- der osmotische Druck.

\begin{werkzeugkasten}[Van-'t-Hoff-Gesetz der Osmose]
\[
  \Pi = c\,R\,T
\]
Mit $c$ = molare Konzentration (mol/L) der gelösten Teilchen, $R$ =
Gaskonstante, $T$ = absolute Temperatur. Für Nicht-Elektrolyte zählt die
Stoffmenge direkt.
\textit{Skript:} Mortimer Kap. 13 (Osmose, osmotischer Druck).
\end{werkzeugkasten}

\begin{achtung}[Darauf achten]
\begin{itemize}
  \item $T$ in Kelvin umrechnen ($25\,°C = 298\ K$).
  \item Konzentration $c = n/V$ mit $V = 100\ mL = 0.100\ L$.
  \item $M(\text{Glukose}) \approx 180.16\ g\,mol^{-1}$.
  \item Einheiten von $R$ zur Einheit von $p$ wählen
        ($R = 0.08314\ L\,bar\,mol^{-1}K^{-1}$ $\Rightarrow$ $\Pi$ in bar).
\end{itemize}
\end{achtung}
